\chapter{Образы}

Образ~--- фундамент всего в Docker. Контейнеры создаются из образов,
реестры хранят образы, Dockerfile описывает, как образ собрать.
В этой главе разобрано, из чего состоит образ, как устроены имена и теги,
почему слои важны для скорости и как безопасно работать с локальным
хранилищем.

\section{Из чего состоит образ}

Образ~--- это не один файл, а набор взаимосвязанных объектов:

\begin{itemize}
  \item \textbf{манифест} (\emph{manifest})~--- JSON-документ со списком слоёв,
        их контрольными суммами и ссылкой на конфигурацию;
  \item \textbf{конфигурация}~--- JSON с метаданными: команда запуска,
        переменные окружения, рабочий каталог, пользователь, порты;
  \item \textbf{слои}~--- tar-архивы с изменениями файловой системы.
\end{itemize}

\begin{figure}[htbp]
\centering
\begin{tikzpicture}[font=\sffamily\scriptsize, x=1cm, y=1cm]
  \draw[fill=warm!8, draw=warm, dashed, rounded corners=2pt]
    (0,4.25) rectangle (9.2,5.1);
  \node[anchor=west, text=warm] at (0.3,4.675)
    {записываемый слой контейнера: изменения исчезают при \cmd{docker rm}};

  \draw[fill=accent!16, draw=accent!55!black, rounded corners=2pt]
    (0,3.15) rectangle (9.2,4.0);
  \node[anchor=west, text=accent!85!black] at (0.3,3.575)
    {слой 3: \cmd{COPY . .}~--- код приложения};
  \node[anchor=east, text=muted] at (8.9,3.575) {1,2 МБ};

  \draw[fill=accent!11, draw=accent!55!black, rounded corners=2pt]
    (0,2.05) rectangle (9.2,2.9);
  \node[anchor=west, text=accent!85!black] at (0.3,2.475)
    {слой 2: \cmd{RUN npm ci}~--- зависимости};
  \node[anchor=east, text=muted] at (8.9,2.475) {85 МБ};

  \draw[fill=accent!7, draw=accent!55!black, rounded corners=2pt]
    (0,0.95) rectangle (9.2,1.8);
  \node[anchor=west, text=accent!85!black] at (0.3,1.375)
    {слой 1: \cmd{FROM node:22-alpine}~--- базовый образ};
  \node[anchor=east, text=muted] at (8.9,1.375) {130 МБ};

  \draw[decorate, decoration={brace, amplitude=5pt, mirror}, draw=accent!60]
    (9.45,0.9) -- (9.45,4.05)
    node[midway, right=6pt, align=left, text=accent!80!black]
      {слои образа\\только для чтения};

  \node[anchor=north west, text=muted, align=left] at (0,-0.15)
    {Манифест перечисляет слои и их digest; конфигурация хранит
     \cmd{CMD}, \cmd{ENV}, \cmd{EXPOSE} и другие метаданные.};
\end{tikzpicture}
\caption{Образ~--- стек неизменяемых слоёв и метаданные; контейнер получает
сверху собственный записываемый слой}\label{fig:layers}
\end{figure}

Слои адресуются по содержимому: \cmd{sha256:...} вычисляется от архива.
Если два образа содержат одинаковый слой, он хранится один раз
и скачивается один раз. Поэтому \cmd{docker pull} второго образа
на той же базе обычно занимает секунды.

\begin{notebox}
Исторически образы Docker состояли из слоёв с файлом \file{layer.json}
и родительскими ссылками. Современные образы следуют спецификации OCI
Image Format: манифест, конфигурация и слои~--- отдельные объекты
в хранилище. Docker, Podman и Kubernetes читают этот формат одинаково.
\end{notebox}

\section{Имена, теги и digest}

Полное имя образа выглядит так:

\begin{center}
\cmd{[registry/][namespace/]repository[:tag][@digest]}
\end{center}

\begin{itemize}
  \item \cmd{registry}~--- адрес реестра. Если не указан, используется
        \cmd{docker.io} (Docker Hub);
  \item \cmd{namespace}~--- владелец или организация. Официальные образы
        живут в пространстве \cmd{library}, поэтому \cmd{nginx}~--- это
        сокращение от \cmd{docker.io/library/nginx};
  \item \cmd{repository}~--- имя образа;
  \item \cmd{tag}~--- человекочитаемая метка версии, по умолчанию
        \cmd{latest};
  \item \cmd{@digest}~--- криптографический хеш манифеста, неизменяемая
        ссылка на конкретную сборку.
\end{itemize}

\begin{code}[language=bash]
$ docker pull nginx
Using default tag: latest
...
$ docker pull nginx:1.27-alpine
$ docker pull nginx@sha256:0f2ab24c6aba...
\end{code}

\begin{warnbox}
\cmd{latest}~--- просто тег по умолчанию, а не \enquote{самая свежая версия}.
Владелец образа может перезаписать его в любой момент. Сборка, которая
год назад использовала \cmd{nginx:latest}, сегодня может получить совсем
другой код. Для воспроизводимости фиксируйте версию (\cmd{nginx:1.27-alpine})
или digest.
\end{warnbox}

Тег~--- это удобная, но подвижная ссылка. Digest~--- точная. Хорошая
практика: в разработке использовать тег, в production~--- тег плюс digest
или неизменяемые теги вида \cmd{1.27.3-alpine}.

\begin{tipbox}
Узнать digest уже скачанного образа:
\cmd{docker image inspect --format '\{\{index .RepoDigests 0\}\}' nginx:1.27-alpine}.
Так можно зафиксировать образ в манифесте поставки, не выискивая хеш
на сайте реестра.
\end{tipbox}

\section{Слои на практике}

Посмотрим на реальный образ.

\begin{code}[language=bash]
$ docker pull nginx:1.27-alpine
$ docker image ls nginx
REPOSITORY   TAG           IMAGE ID       CREATED       SIZE
nginx        1.27-alpine   a8b1c3d4e5f6   2 weeks ago   48.2MB
$ docker history nginx:1.27-alpine
IMAGE          CREATED       CREATED BY                                      SIZE
a8b1c3d4e5f6   2 weeks ago   CMD ["nginx" "-g" "daemon off;"]                0B
<missing>      2 weeks ago   STOPSIGNAL SIGQUIT                              0B
<missing>      2 weeks ago   EXPOSE map[80/tcp:{}]                           0B
...
<missing>      2 weeks ago   /bin/sh -c #(nop) ADD file:... in /             7.8MB
\end{code}

Обратите внимание: инструкции, которые меняют только метаданные
(\cmd{CMD}, \cmd{EXPOSE}, \cmd{STOPSIGNAL}), имеют размер \cmd{0B}
и не создают новых слоёв. Слои появляются там, где меняется файловая
система: \cmd{RUN}, \cmd{COPY}, \cmd{ADD}.

Команда \cmd{docker history} показывает и \cmd{<missing>}~--- это
нормально: промежуточные образы старых сборок не хранятся, но их слои
остаются частью итогового образа.

\section{Основные команды работы с образами}

\begin{center}
\small
\begin{tabularx}{\textwidth}{@{}lX@{}}
\toprule
\textbf{Команда} & \textbf{Назначение} \\
\midrule
\cmd{docker pull IMAGE} & скачать образ из реестра \\
\cmd{docker images} & список локальных образов (алиас \cmd{docker image ls}) \\
\cmd{docker image inspect IMAGE} & метаданные в JSON \\
\cmd{docker history IMAGE} & слои и создавшие их инструкции \\
\cmd{docker tag SRC DST} & добавить ещё одно имя (тег) к образу \\
\cmd{docker rmi IMAGE} & удалить образ \\
\cmd{docker image prune} & удалить \enquote{висячие} образы \\
\cmd{docker save -o file.tar IMAGE} & сохранить образ в архив \\
\cmd{docker load -i file.tar} & загрузить образ из архива \\
\cmd{docker export CONTAINER} & экспортировать файловую систему контейнера \\
\bottomrule
\end{tabularx}
\end{center}

\begin{code}[language=bash]
$ docker image inspect nginx:1.27-alpine \
    --format '{{.Config.Cmd}}'
[/docker-entrypoint.sh nginx -g daemon off;]
$ docker image inspect nginx:1.27-alpine \
    --format '{{json .Config.Env}}'
["PATH=/usr/local/sbin:...","NGINX_VERSION=1.27.3",...]
$ docker tag nginx:1.27-alpine registry.example.com/team/nginx:1.27.3
$ docker rmi nginx:1.27-alpine
\end{code}

\begin{notebox}
\cmd{docker rmi} удаляет только имя. Слои освободятся, когда на образ
не останется ссылок: тегов и контейнеров (в том числе остановленных).
Если образ используется контейнером, команда откажется его удалять,
пока вы не удалите контейнер или не добавите \flag{-f}.
\end{notebox}

\begin{warnbox}
\cmd{docker image prune -a} удаляет все образы, не используемые
контейнерами, включая те, что вы скачали для работы завтра.
Не выполняйте эту команду \enquote{на всякий случай} на рабочей машине:
повторное скачивание может быть долгим, а локально собранные образы
исчезнут безвозвратно.
\end{warnbox}

\section{\cmd{save} и \cmd{export}: не одно и то же}

Две команды легко перепутать.

\begin{itemize}
  \item \cmd{docker save} сохраняет \emph{образ} целиком: все слои
        и метаданные. \cmd{docker load} восстанавливает его один в один.
        Это способ перенести образ в изолированную сеть.
  \item \cmd{docker export} сохраняет \emph{файловую систему контейнера}
        одним слоем, без истории и метаданных образа. Результат можно
        использовать как rootfs, но не как полноценный образ с
        \cmd{CMD}, \cmd{ENV} и \cmd{ENTRYPOINT}.
\end{itemize}

\begin{code}[language=bash]
$ docker save -o nginx.tar nginx:1.27-alpine
$ docker load -i nginx.tar
$ docker export web > web-rootfs.tar
\end{code}

\begin{tipbox}
Если нужно передать образ между машинами без реестра, используйте
\cmd{docker save} вместе со сжатием:
\cmd{docker save nginx:1.27-alpine | gzip > nginx.tar.gz}.
Обратная операция~--- \cmd{gunzip -c nginx.tar.gz | docker load}.
\end{tipbox}

\section{Как читать размер образа}

\cmd{docker image ls} показывает размер распакованных слоёв, приписывая
каждый общий слой первому образу, который его использует. Поэтому сумма
размеров в списке может не совпадать с занятым диском.
Реальную картину даёт \cmd{docker system df -v}: он показывает,
сколько места занимают образы, контейнеры, тома и кэш сборки.

\begin{code}[language=bash]
$ docker system df
TYPE            TOTAL     ACTIVE    SIZE      RECLAIMABLE
Images          12        3         2.41GB    1.82GB (75%)
Containers      5         2         12.3MB    8.1MB (65%)
Local Volumes   4         2         210MB     96MB (45%)
Build Cache     34        0         1.1GB     1.1GB
\end{code}

\begin{notebox}
Большой образ не обязательно означает медленный запуск: контейнер
не копирует слои, а использует их через overlayfs. Но большой образ
дольше скачивать и дольше распаковывать при первом запуске на новой
машине. Как уменьшать размеры~--- в главе~13.
\end{notebox}

\section{Образы и архитектуры}

Один тег может указывать на несколько сборок для разных архитектур.
Например, \cmd{nginx:alpine} для \cmd{amd64} и \cmd{arm64}~--- разные
манифесты, объединённые списком (manifest list). При скачивании Docker
автоматически выбирает подходящий вариант.

\begin{code}[language=bash]
$ docker manifest inspect nginx:1.27-alpine
{
   "manifests": [
      { "platform": { "architecture": "amd64", "os": "linux" }, ... },
      { "platform": { "architecture": "arm64", "os": "linux" }, ... }
   ]
}
\end{code}

Если образа для вашей архитектуры нет, вы получите предупреждение
о несовпадении платформы или ошибку \cmd{exec format error}. Решения:
использовать мультиархитектурный образ, собрать свой (раздел~10.5)
или включить эмуляцию через QEMU (медленно, только для отладки).

\begin{summarybox}
\begin{itemize}
  \item Образ~--- манифест, конфигурация и набор неизменяемых слоёв.
  \item Имя образа включает реестр, пространство имён, репозиторий, тег
        и, при необходимости, digest.
  \item Тег подвижен, digest неизменяем; для воспроизводимости фиксируйте
        версию.
  \item Слои разделяются между образами и кэшируются; метаданные новых
        слоёв не создают.
  \item \cmd{save} переносит образ, \cmd{export}~--- только файловую
        систему контейнера.
\end{itemize}
\end{summarybox}

\section*{Вопросы для самопроверки}
\begin{selfcheck}
  \item Из каких трёх частей состоит образ?
  \item Почему \cmd{nginx} и \cmd{docker.io/library/nginx:latest}~---
        одно и то же?
  \item Чем тег отличается от digest и почему это важно для production?
  \item Какие инструкции Dockerfile создают слои, а какие~--- нет?
  \item Чем \cmd{docker save} отличается от \cmd{docker export}?
  \item Почему сумма размеров в \cmd{docker images} может превышать
        реально занятое место?
\end{selfcheck}
